\section{FORMAL RESTATEMENT}
\textbf{Erd\H{o}s \#223; reconstruction of the asymptotic diameter bound, 2026-09-23.}
For $d\geq2$ and $n\geq2$, let $f_d(n)$ be the maximum number of unordered pairs at distance one among $n$ distinct points in $\mathbb R^d$ whose diameter is one. Estimate $f_d(n)$.

\section{QUICK LITERATURE AND CONTEXT CHECK}
The known low-dimensional answers are $f_2(n)=n$ for $n\geq3$ and $f_3(n)=2n-2$ for $n\geq4$. The small exceptions are $f_d(2)=1$ and $f_d(3)=3$. For each fixed $d\geq4$, writing $p=\lfloor d/2\rfloor$, the leading asymptotic is $(p-1)n^2/(2p)$. Swanepoel subsequently determined the exact values for all sufficiently large $n$ depending on $d$. This file reconstructs the leading asymptotic and sharp lower-dimensional examples; it uses the established low-dimensional upper bounds as named external inputs rather than reproving them.

\section{OBJECTIVE AND ATTACK PLAN}
For the lower bound, use short arcs on circles in orthogonal planes so that every cross-pair is a diameter and no within-circle distance is too large. For the upper bound, show that a certain complete multipartite graph cannot be realized even as a unit-distance graph, and invoke the Erd\H{o}s--Stone theorem.

\section{WORK}
\textbf{The higher-dimensional construction.} For $d\geq4$ let $p=\lfloor d/2\rfloor\geq2$, and choose $p$ mutually orthogonal two-dimensional linear subspaces. In each take a circle centred at zero of radius $1/\sqrt2$. Choose $n_i$ distinct points on an arc of angular length at most $\pi/2$, with $\sum_i n_i=n$ and all $n_i$ differing by at most one. Two points on different circles have squared distance $1/2+1/2=1$. Two points on the same arc have distance at most $2(1/\sqrt2)\sin(\pi/4)=1$. For all sufficiently large $n$ the set therefore has diameter one and at least
\[
\sum_{i<j}n_i n_j
=\frac12\left(n^2-\sum_i n_i^2\right)
=\frac{p-1}{2p}n^2-O_p(1)
\tag{1}
\]
diameter pairs. The short-arc restriction is essential for a diameter problem; arbitrary points on a whole circle could be farther than one apart.

\textbf{A forbidden multipartite graph.} A unit-distance graph in $\mathbb R^d$ cannot contain $K_{p+1}(3)$, the complete $(p+1)$-partite graph with three vertices in each part. Suppose otherwise, with parts $A_1,\ldots,A_{p+1}$. Let $V_i$ be the linear space spanned by differences of points of $A_i$. It has dimension at least two. Indeed any point from another part is at distance one from all three points of $A_i$, and a line intersects a sphere in at most two points; the three points therefore cannot be collinear.

For $a,a'\in A_i$ and $b,b'\in A_j$, $i\ne j$, all four cross-distances are one. Expanding their squared norms and subtracting gives
\[
|a-b|^2-|a-b'|^2-|a'-b|^2+|a'-b'|^2
=-2(a-a')\cdot(b-b')=0.
\]
Hence $V_i\perp V_j$. Mutually orthogonal subspaces have direct sum, so
\[
d\geq\sum_{i=1}^{p+1}\dim V_i\geq2(p+1)>d,
\]
a contradiction.

\textbf{Declared external graph theorem.} Erd\H{o}s--Stone says that for a fixed graph $F$ of chromatic number $r\geq3$,
\[
\operatorname{ex}(n,F)
=\left(1-\frac1{r-1}+o(1)\right)\frac{n^2}{2}.
\tag{2}
\]
Apply (2) to $F=K_{p+1}(3)$, whose chromatic number is $p+1$. The diameter graph of any candidate configuration is a unit-distance graph, and the preceding argument excludes this $F$ as a subgraph, not merely as an induced subgraph. Therefore
\[
f_d(n)\leq\frac{p-1}{2p}n^2+o_d(n^2).
\tag{3}
\]
Combining (1) and (3) proves the asserted asymptotic for every fixed $d\geq4$.

\textbf{Low-dimensional sharp examples.} In the plane take a point $O$ and $n-1$ distinct points on a unit-circle arc about $O$ whose endpoints make angle $\pi/3$, including both endpoints. For $n\geq3$, all $n-1$ radial pairs have distance one, as does the endpoint pair; every other pair on the arc has smaller distance. This gives exactly $n$ diameter pairs.

In three dimensions take $A=(-1/2,0,0)$ and $B=(1/2,0,0)$. The common intersection of their unit spheres is the circle $x=0$, $y^2+z^2=3/4$. For $n\geq4$ choose $n-2$ distinct points on an arc of this circle whose endpoints have chord length one, including both endpoints. Its angular length is $2\arcsin(1/\sqrt3)<\pi$, so all within-arc distances are at most one, with equality only for the endpoints. All $2(n-2)$ cross-pairs to $A,B$ have distance one, as do $AB$ and the endpoint pair, giving exactly $2n-2$ diameter pairs.

The matching external upper bounds are the planar diameter theorem of Hopf--Pannwitz and the V\'azsonyi theorem proved independently by Gr\"unbaum, Heppes, and Straszewicz. Together with these explicit examples they give the exact low-dimensional formulas stated above. For $n=2,3$, a pair and an equilateral triangle attain the trivial upper bound $\binom n2$.

\section{VERIFICATION AND SCOPE}
The asymptotic is for fixed $d$ as $n\to\infty$. Its geometric obstruction and lower construction are proved here; Erd\H{o}s--Stone and the two low-dimensional upper theorems are declared dependencies. No claim is made to derive Swanepoel's sharper exact eventual formulas or all extremal configurations. The requested leading-order estimate and the exact two- and three-dimensional values follow from the stated inputs, with the small-order exceptions handled explicitly.

\section{SOURCES}
\begin{itemize}
\item Current statement and original attributions: \url{https://www.erdosproblems.com/223}.
\item K. J. Swanepoel, \emph{Unit distances and diameters in Euclidean spaces}, Sections 1.1--1.2 and the references there for the original low-dimensional results; the paper also gives the stronger exact eventual results: \url{https://arxiv.org/abs/0707.0213}.
\end{itemize}
\section{CONCLUSION}
\textbf{Attempt status: solved at the stated estimation scope using named external theorems.} The higher-dimensional leading asymptotic and the exact low-dimensional values follow. This is a reconstruction of known results, with no novelty claim.
\textbf{Completion rate estimate: 100\%.}
